% Options for packages loaded elsewhere
% Options for packages loaded elsewhere
\PassOptionsToPackage{unicode}{hyperref}
\PassOptionsToPackage{hyphens}{url}
%
\documentclass[
  11pt,
  ignorenonframetext,
  aspectratio=169,
  11pt]{beamer}
\newif\ifbibliography
\usepackage{pgfpages}
\setbeamertemplate{caption}[numbered]
\setbeamertemplate{caption label separator}{: }
\setbeamercolor{caption name}{fg=normal text.fg}
\beamertemplatenavigationsymbolsempty
% remove section numbering
\setbeamertemplate{part page}{
  \centering
  \begin{beamercolorbox}[sep=16pt,center]{part title}
    \usebeamerfont{part title}\insertpart\par
  \end{beamercolorbox}
}
\setbeamertemplate{section page}{
  \centering
  \begin{beamercolorbox}[sep=12pt,center]{section title}
    \usebeamerfont{section title}\insertsection\par
  \end{beamercolorbox}
}
\setbeamertemplate{subsection page}{
  \centering
  \begin{beamercolorbox}[sep=8pt,center]{subsection title}
    \usebeamerfont{subsection title}\insertsubsection\par
  \end{beamercolorbox}
}
% Prevent slide breaks in the middle of a paragraph
\widowpenalties 1 10000
\raggedbottom
\AtBeginPart{
  \frame{\partpage}
}
\AtBeginSection{
  \ifbibliography
  \else
    \frame{\sectionpage}
  \fi
}
\AtBeginSubsection{
  \frame{\subsectionpage}
}
\usepackage{iftex}
\ifPDFTeX
  \usepackage[T1]{fontenc}
  \usepackage[utf8]{inputenc}
  \usepackage{textcomp} % provide euro and other symbols
\else % if luatex or xetex
  \usepackage{unicode-math} % this also loads fontspec
  \defaultfontfeatures{Scale=MatchLowercase}
  \defaultfontfeatures[\rmfamily]{Ligatures=TeX,Scale=1}
\fi
\usepackage{lmodern}

\usetheme[]{Madrid}
\ifPDFTeX\else
  % xetex/luatex font selection
\fi
% Use upquote if available, for straight quotes in verbatim environments
\IfFileExists{upquote.sty}{\usepackage{upquote}}{}
\IfFileExists{microtype.sty}{% use microtype if available
  \usepackage[]{microtype}
  \UseMicrotypeSet[protrusion]{basicmath} % disable protrusion for tt fonts
}{}
% Map common Unicode glyphs to LaTeX commands to avoid missing glyph errors
\IfFileExists{newunicodechar.sty}{\usepackage{newunicodechar}}{\usepackage{newunicodechar}}
\newunicodechar{•}{\textbullet}
\newunicodechar{−}{\textminus}
\newunicodechar{⋅}{\ensuremath{\cdot}}
\newunicodechar{𝛿}{\ensuremath{\delta}}
\makeatletter
\@ifundefined{KOMAClassName}{% if non-KOMA class
  \IfFileExists{parskip.sty}{%
    \usepackage{parskip}
  }{% else
    \setlength{\parindent}{0pt}
    \setlength{\parskip}{6pt plus 2pt minus 1pt}}
}{% if KOMA class
  \KOMAoptions{parskip=half}}
\makeatother


\usepackage{longtable,booktabs,array}
\newcounter{none} % for unnumbered tables
\usepackage{calc} % for calculating minipage widths
\usepackage{caption}
% Make caption package work with longtable
\makeatletter
\def\fnum@table{\tablename~\thetable}
\makeatother
\usepackage{graphicx}
\makeatletter
\newsavebox\pandoc@box
\newcommand*\pandocbounded[1]{% scales image to fit in text height/width
  \sbox\pandoc@box{#1}%
  \Gscale@div\@tempa{\textheight}{\dimexpr\ht\pandoc@box+\dp\pandoc@box\relax}%
  \Gscale@div\@tempb{\linewidth}{\wd\pandoc@box}%
  \ifdim\@tempb\p@<\@tempa\p@\let\@tempa\@tempb\fi% select the smaller of both
  \ifdim\@tempa\p@<\p@\scalebox{\@tempa}{\usebox\pandoc@box}%
  \else\usebox{\pandoc@box}%
  \fi%
}
% Set default figure placement to htbp
\def\fps@figure{htbp}
\makeatother





\setlength{\emergencystretch}{3em} % prevent overfull lines

\providecommand{\tightlist}{%
  \setlength{\itemsep}{0pt}\setlength{\parskip}{0pt}}



 


% Basic font and theme packages
\usepackage[T1]{fontenc}
\usepackage[utf8]{inputenc}
\usepackage{mathpazo}
\usepackage[scaled=0.92]{helvet}
\usepackage{courier}
\usecolortheme{default}
\usefonttheme{serif}

% Common packages
\usepackage{booktabs}
\usepackage{array}
\usepackage{ragged2e}
\usepackage{tikz}
\usepackage{amsmath}
\usepackage{amssymb}
\usepackage{microtype}

% Davidson-style academic palette and spacing
\definecolor{DavidsonBlue}{HTML}{1F4E79}
\definecolor{DavidsonBlueDark}{HTML}{173A5E}
\definecolor{SoftGray}{HTML}{F4F6F8}
\definecolor{AccentGold}{HTML}{B8860B}
\definecolor{TextGray}{HTML}{333333}

\setbeamercolor{title}{fg=white,bg=DavidsonBlue}
\setbeamercolor{frametitle}{fg=white,bg=DavidsonBlue}
\setbeamercolor{structure}{fg=DavidsonBlueDark}
\setbeamercolor{normal text}{fg=TextGray,bg=white}
\setbeamercolor{block title}{fg=white,bg=DavidsonBlueDark}
\setbeamercolor{block body}{fg=TextGray,bg=SoftGray}
\setbeamercolor{item projected}{fg=white,bg=DavidsonBlue}
\setbeamercolor{footline}{fg=white,bg=DavidsonBlueDark}
\setbeamercolor{author in head/foot}{fg=white,bg=DavidsonBlueDark}
\setbeamercolor{title in head/foot}{fg=white,bg=DavidsonBlueDark}
\setbeamercolor{date in head/foot}{fg=white,bg=DavidsonBlueDark}

\setbeamertemplate{navigation symbols}{}
\setbeamertemplate{items}[circle]
\setbeamertemplate{blocks}[rounded][shadow=false]

\setbeamerfont{title}{series=\bfseries,size=\Large}
\setbeamerfont{subtitle}{size=\large}
\setbeamerfont{frametitle}{series=\bfseries,size=\large}
\setbeamerfont{author}{size=\normalsize}
\setbeamerfont{date}{size=\small}

\setbeamertemplate{frametitle}{%
  \nointerlineskip
  \begin{beamercolorbox}[wd=\paperwidth,ht=3.1ex,dp=1.1ex,leftskip=0.45cm,rightskip=0.2cm]{frametitle}
    \usebeamerfont{frametitle}\insertframetitle
  \end{beamercolorbox}%
}

\setbeamertemplate{footline}{%
  \leavevmode%
  \hbox{%
    \begin{beamercolorbox}[wd=.73\paperwidth,ht=2.6ex,dp=1.1ex,leftskip=0.35cm]{author in head/foot}%
      \usebeamerfont{author in head/foot}Supply-Side Teacher Forecasting Framework
    \end{beamercolorbox}%
    \begin{beamercolorbox}[wd=.19\paperwidth,ht=2.6ex,dp=1.1ex,center]{title in head/foot}%
      Tanzania
    \end{beamercolorbox}%
    \begin{beamercolorbox}[wd=.08\paperwidth,ht=2.6ex,dp=1.1ex,rightskip=0.25cm]{date in head/foot}%
      \hfill\insertframenumber/\inserttotalframenumber
    \end{beamercolorbox}%
  }
}

\newcommand{\highlight}[1]{\textcolor{DavidsonBlueDark}{\textbf{#1}}}
\newcommand{\gold}[1]{\textcolor{AccentGold}{\textbf{#1}}}
\makeatletter
\@ifpackageloaded{caption}{}{\usepackage{caption}}
\AtBeginDocument{%
\ifdefined\contentsname
  \renewcommand*\contentsname{Table of contents}
\else
  \newcommand\contentsname{Table of contents}
\fi
\ifdefined\listfigurename
  \renewcommand*\listfigurename{List of Figures}
\else
  \newcommand\listfigurename{List of Figures}
\fi
\ifdefined\listtablename
  \renewcommand*\listtablename{List of Tables}
\else
  \newcommand\listtablename{List of Tables}
\fi
\ifdefined\figurename
  \renewcommand*\figurename{Figure}
\else
  \newcommand\figurename{Figure}
\fi
\ifdefined\tablename
  \renewcommand*\tablename{Table}
\else
  \newcommand\tablename{Table}
\fi
}
\@ifpackageloaded{float}{}{\usepackage{float}}
\floatstyle{ruled}
\@ifundefined{c@chapter}{\newfloat{codelisting}{h}{lop}}{\newfloat{codelisting}{h}{lop}[chapter]}
\floatname{codelisting}{Listing}
\newcommand*\listoflistings{\listof{codelisting}{List of Listings}}
\makeatother
\makeatletter
\makeatother
\makeatletter
\@ifpackageloaded{caption}{}{\usepackage{caption}}
\@ifpackageloaded{subcaption}{}{\usepackage{subcaption}}
\makeatother

\usepackage{bookmark}
\IfFileExists{xurl.sty}{\usepackage{xurl}}{} % add URL line breaks if available
\urlstyle{same}
\hypersetup{
  hidelinks,
  pdfcreator={LaTeX via pandoc}}


\author{}
\date{}

\begin{document}

\begin{frame}[plain]
  \begin{tikzpicture}[remember picture,overlay]
    \fill[DavidsonBlue] (current page.south west) rectangle (current page.north east);
    \fill[white, opacity=0.06] ([xshift=-1cm,yshift=-0.5cm]current page.north east) circle (2.5cm);
    \fill[white, opacity=0.05] ([xshift=1.5cm,yshift=0.8cm]current page.south west) circle (3cm);
  \end{tikzpicture}
  \vspace{0.9cm}
  {\usebeamerfont{title}\color{white}Supply-Side Teacher Forecasting Framework\par}
  \vspace{0.2cm}
  {\usebeamerfont{subtitle}\color{white!92}Primary and Secondary Education in Tanzania\par}
  \vspace{0.7cm}
  {\color{white}\large Prepared for: \par}
  \vspace{0.15cm}
  {\color{white!92}\normalsize Ministry of Education, Science and Technology (MOEST)\\
  President's Office - Regional Administration and Local Government (PO-RALG)\par}
  \vfill
  {\color{white!85}\small A technical framework for workforce planning, budgeting, and teacher deployment.\par}
  \vspace{0.35cm}
    {\color{white!80}\small March 2026\par}
\end{frame}

\begin{frame}{Strategic Context}
\begin{columns}[T,onlytextwidth]
\column{0.58\textwidth}
\begin{itemize}
  \item Tanzania has expanded access to basic and secondary education through fee-free education, curriculum reforms, and school expansion.
  \item These gains have increased \highlight{demand for teachers}, especially across growing enrolment areas.
  \item Persistent constraints include fiscal pressures, uneven deployment, limited training capacity in some subjects, and attrition.
  \item Acute shortages remain visible in \gold{STEM}, languages, special needs education, and remote localities.
\end{itemize}
\column{0.38\textwidth}
\begin{block}{Why forecasting?}
\justifying
A regular forecasting system enables MOEST and PO-RALG to shift from reactive staffing decisions to \highlight{proactive workforce planning}.
\end{block}
\end{columns}
\end{frame}

\begin{frame}{Purpose and Scope}
\begin{block}{Purpose}
\justifying
To provide a \highlight{clear, repeatable, and transparent} process for forecasting the supply of primary and secondary teachers in Tanzania.
\end{block}
\vspace{0.15cm}
\begin{columns}[T,onlytextwidth]
\column{0.52\textwidth}
\textbf{What the framework supports}
\begin{itemize}
  \item Recruitment planning
  \item Budgeting and wage-bill discussions
  \item Teacher deployment and post allocation
  \item Scholarship and training decisions
\end{itemize}
\column{0.46\textwidth}
\textbf{What it covers}
\begin{itemize}
  \item Current stock of teachers
  \item New graduates and re-entrants
  \item Retirements, mortality, and attrition
  \item Disaggregation by level, subject, gender, and region
\end{itemize}
\end{columns}
\end{frame}

\begin{frame}{Institutional Arrangements and Governance}
\begin{itemize}
  \item \highlight{MOEST}: policy direction, national standards, technical oversight, and use of forecasts for teacher education and sector planning.
  \item \highlight{PO-RALG}: deployment, LGA coordination, stock and attrition data, and local interpretation of forecast results.
  \item \highlight{Teacher education institutions, TCU, NACTVET}: intake, completion, and subject specialization data.
  \item \highlight{Ministry of Finance, NBS, development partners}: wage-bill implications, demographic projections, and technical support.
\end{itemize}
\vspace{0.25cm}
\begin{block}{Recommended coordination mechanism}
A joint \highlight{Teacher Workforce Planning Technical Working Group} should review data, validate assumptions, agree scenarios, and sign off forecast outputs.
\end{block}
\end{frame}

\begin{frame}{Forecasting Methodology}
\begin{columns}[T,onlytextwidth]
\column{0.56\textwidth}
\textbf{Core approach}
\begin{itemize}
  \item Stock-and-flow model as the main forecasting backbone
  \item Trend-based and ratio-based projections where detailed data are incomplete
  \item Scenario analysis to test plausible alternative futures
  \item Structured expert judgment to refine assumptions
\end{itemize}
\column{0.40\textwidth}
\begin{block}{Why this approach?}
\justifying
It is \highlight{transparent}, uses relatively low data requirements, supports \highlight{data-driven decisions}, and can be institutionalized while maintaining attention to \highlight{equity, gender, and inclusion}.
\end{block}
\end{columns}
\end{frame}

\begin{frame}{Core Stock-Flow Identity}
\begin{block}{Dynamic accounting relationship}
\[
\mathrm{Stock}_{t+1} = \mathrm{Stock}_{t} + \mathrm{Inflows}_{t} - \mathrm{Outflows}_{t}
\]
\end{block}
\vspace{0.15cm}
\begin{columns}[T,onlytextwidth]
\column{0.32\textwidth}
\textbf{Stock}
\begin{itemize}
  \item Teachers in post
  \item Actively available to teach
\end{itemize}
\column{0.32\textwidth}
\textbf{Inflows}
\begin{itemize}
  \item New graduates recruited
  \item Re-entries
  \item Transfers in
\end{itemize}
\column{0.32\textwidth}
\textbf{Outflows}
\begin{itemize}
  \item Retirements
  \item Mortality
  \item Attrition and exits
\end{itemize}
\end{columns}
\vspace{0.2cm}
\begin{block}{Interpretation}
The model updates the teacher workforce \highlight{year by year}, making it suitable for annual planning and budgeting cycles.
\end{block}
\end{frame}

\begin{frame}{Inflows and Outflows}
\begin{columns}[T,onlytextwidth]
\column{0.49\textwidth}
\begin{block}{Inflows}
\[
H_t = \mathrm{Intake}_{t-k}\cdot c \cdot \delta
\]
\begin{itemize}
  \item Training intake drives future graduate supply
  \item $c$: completion rate
  \item $\delta$: recruitment or absorption rate
  \item Re-entries add a smaller but relevant source of supply
\end{itemize}
\end{block}
\column{0.49\textwidth}
\begin{block}{Outflows}
\[
\mathrm{Outflow}_t = \mathrm{Ret}_t + \mathrm{Mort}_t + \mathrm{Attr}_t
\]
\begin{itemize}
  \item Retirement is relatively predictable
  \item Mortality is usually small but non-negligible
  \item Attrition is the \gold{most policy-sensitive} component
\end{itemize}
\end{block}
\end{columns}
\end{frame}

\begin{frame}{Scenarios and Validation}
\begin{columns}[T,onlytextwidth]
\column{0.47\textwidth}
\textbf{Scenario design}
\begin{itemize}
  \item Baseline: continuation of current trends
  \item Optimistic: stronger training capacity and retention
  \item Pessimistic: fiscal tightening or higher attrition
\end{itemize}
\column{0.49\textwidth}
\textbf{Validation}
\begin{itemize}
  \item Back-testing with historical data
  \item Mean Absolute Error (MAE)
  \item Root Mean Square Error (RMSE)
  \item Mean Absolute Percentage Error (MAPE)
\end{itemize}
\end{columns}
\vspace{0.2cm}
\begin{block}{Planning value}
Scenario analysis supports \highlight{risk-aware teacher workforce planning} rather than dependence on a single point forecast.
\end{block}
\end{frame}

\begin{frame}{Implementation Roadmap}
\begin{enumerate}
  \item \highlight{Phase 1: Pilot and prepare}
  \begin{itemize}
    \item Map data sources, assess quality, integrate historical data
    \item Pilot a basic stock-flow model for selected regions or subjects
  \end{itemize}
  \item \highlight{Phase 2: Scale up}
  \begin{itemize}
    \item Extend to all regions and subject groups
    \item Align forecasting cycle with budgeting and recruitment decisions
  \end{itemize}
  \item \highlight{Phase 3: Consolidate and enhance}
  \begin{itemize}
    \item Introduce more advanced statistical modules where data permit
    \item Integrate supply and demand forecasting for fuller workforce planning
  \end{itemize}
\end{enumerate}
\end{frame}

\begin{frame}{Conclusion and Policy Message}
\begin{block}{Core message}
\justifying
A practical supply-side forecasting framework allows Tanzania's education system to move from \highlight{reactive staffing responses} toward \highlight{strategic teacher workforce management}.
\end{block}
\vspace{0.15cm}
\begin{itemize}
  \item It strengthens recruitment planning, deployment, and budgeting.
  \item It makes assumptions explicit and testable.
  \item It improves attention to equity, gender balance, and underserved areas.
  \item It can begin with existing tools and evolve as data quality improves.
\end{itemize}
\vspace{0.2cm}
\centering
\Large\gold{Every learner should have access to a qualified teacher.}
\end{frame}




\end{document}
